\section{Introduction}\label{sec:intro}

Quantitative finance runs on licensed data: vendor feeds, index
histories, alternative datasets. Licenses end. When one does, the
data must stop being used, including \emph{inside trained models}.
This is the problem machine unlearning addresses, and it arrives
here with a contractual rather than privacy motivation. The field's central obstacle is that
the ground truth of forgetting, retraining from scratch without the
deleted data, is unaffordable at LLM scale. Approximate methods are
therefore scored against
proxies~\cite{nguyen2025survey}, and the field's
verification tools are contested~\cite{zhang2024mia_cannot}.

Financial forecasters invert the economics: they retrain in minutes.
That affordability buys the study the field cannot run at scale. We
answer every question about unlearning against \emph{exact retraining
oracles}, 3{,}200 of them, on an S\&P~500 volatility-forecasting
panel, with every comparison paired by trading day.

One result reorganizes the rest. Deleting a clean, legitimate year of
market data is not a cost. In the next bear market it is a gain, and
the gain is largest for the models that memorized the year most. That
cuts against the long-tail account of memorization, in which
memorizing rare data is what buys generalization and removing it
should cost accuracy~\cite{feldman2020memorization}. Nothing is
wrong with the year we delete. What changes is the setting: a license
expires, the model is judged on a later regime, and under that shift
the value of a memorized era can invert. The unlearning literature
could not have found this. Approximate unlearning shows utility
falling, never rising. Seeing it rise takes the exact counterfactual,
which at this scale we can afford. Three research questions organize
the paper.

\textbf{RQ1: What does deleting a licensed era mean, and what does
exact deletion do?} \emph{Exact deletion is nearly inert in calm
markets, and under stress its sign is a property of the architecture,
not of the data.} First the question has to be made well-posed.
``Delete year $E$'' is not, because overlapping windows, rolling
features, and pipeline constants touch $E$ beyond its calendar
bounds. We formalize a membership lattice of four increasingly
strict readings and adopt the window-level reading as the
paper-wide default. Measured against oracles, deletion moves a calm
test year by almost nothing. In the next bear market the same
deletion moves an order of magnitude more, and every deployable
architecture gains from removing the crisis year. The two
architectures below a standard volatility benchmark do not gain,
and they are the two no one should deploy. Deletion exposure is
therefore something to screen for, not a fixed cost to budget.
Within a single training window, the effect follows the deleted
era's position. The earliest years are inert, the years just before
the test period cost skill, and deleting the crisis year helps.

\textbf{RQ2: Can approximate unlearning substitute for retraining,
and what is the right target?} \emph{The right target is the
counterfactual, and the counterfactual still predicts the deleted
era.} The only prior forecasting-specific
method~\cite{wang2026tsunlearn} drives forgotten predictions toward
noise. Yet the counterfactual, retrained without the crisis year,
still predicts that year almost as well as the reference, because
nearly all of that skill is generalization rather than memory.
Approaching the noise target therefore surrenders live skill.
Compliance means matching the counterfactual. On substitution the
answer is no. In a single fold several deployable
architecture-method pairs look adequate. Across all four folds,
exactly one brackets the oracle, one more stays within half the
gap, and the rest under-forget or overshoot.

\textbf{RQ3: Can deletion be audited?} \emph{Only when the model
memorized the era. Auditability is a property of what was learned, not
of what was deleted.} An auditor calibrated against a population of
retrainings separates crisis-era deletions easily, but across the
full grid the separation falls to chance, and seventeen per cent of
legitimate deletions leave a trace too faint to verify. A null
audit is therefore not evidence of compliance. Window-level
auditing stays near chance everywhere, and treating overlapping
windows as i.i.d.\ overstates audit confidence. Audit granularity
belongs in the data contract.

\textbf{Contributions.} (i)~The sign result: exact deletion of a
crisis era improves every deployable architecture in all four
market regimes, and the sign is carried by the architecture, not by
the data. (ii)~The instrument: 3{,}200 exact retraining oracles and
the full oracle--method--audit matrix, released for reuse. (iii)~A
membership lattice that makes ``delete year $E$'' well-posed under
overlapping windows, with measurements of what each reading costs
in data and in measured memory. (iv)~The target result: the
counterfactual still predicts the deleted era, so compliance and
privacy forgetting are different objectives, and pursuing the
second surrenders the model. (v)~The audit result: audit strength
tracks what the model memorized, not what was deleted, so
legitimate deletions can be unverifiable at any granularity.
